% TOP_PROBLEM: {"id": 3106, "title": "Shannon capacity of the seven-cycle", "category_id": 3, "category": {"id": 3, "name": "graph_theory", "display_name": "Graph Theory", "description": "Problems involving graphs, networks, and their properties.", "slug": "graph-theory", "order_index": 3, "created_at": "2026-07-31T15:26:25.671Z"}, "status": "open", "problem_number": "OPG-36879", "set_id": 12, "statusQualification": "The conventional Shannon-capacity normalization is used; the source’s extra factor 1/k under the kth root does not change its limiting value."}
% FIRST_POSTED: 2026-10-01
% ATTEMPT_STATUS: unresolved
% UnsolvedMath ID 3106; source catalogue code OPG-36879
% !TeX program = lualatex
% GPT-6 Astra Ultra research attempt, 2026-10-01; self-contained partial work.
% Completion estimate: no numerical percentage assigned; full problem unresolved.
\documentclass[11pt,a4paper]{article}
\usepackage[margin=25mm]{geometry}
\usepackage{fontspec}
\setmainfont{lmroman10-regular.otf}[BoldFont=lmroman10-bold.otf,
 ItalicFont=lmroman10-italic.otf,BoldItalicFont=lmroman10-bolditalic.otf]
\usepackage{amsmath,amssymb,mathtools}
\usepackage{unicode-math}
\setmathfont{latinmodern-math.otf}
\AtBeginDocument{\let\setminus\smallsetminus}
\usepackage{xurl,hyperref}
\hypersetup{colorlinks=true,urlcolor=blue}
\setcounter{secnumdepth}{0}
\setlength{\parindent}{0pt}
\setlength{\parskip}{5pt}
\setlength{\emergencystretch}{3em}
\begin{document}
\section*{Shannon capacity of the seven-cycle}\label{3106}
{\small UnsolvedMath ID 3106 · OPG-36879 · Graph Theory · OpenGarden}
\subsection{Definitions and mathematical statement}
Let $C_7$ be the simple cycle with vertices $\mathbb Z/7\mathbb Z$ and
edges between residues whose difference is $1$ or $-1$ modulo $7$.
For $k\ge1$, its $k$th strong power $C_7^{\boxtimes k}$ has vertex set
$(\mathbb Z/7\mathbb Z)^k$. Distinct tuples $x,y$ are adjacent if,
in every coordinate, they are equal or adjacent in $C_7$.

For a finite graph $H$, the independence number $\alpha(H)$ is the
largest cardinality of a vertex set containing no adjacent pair.
The Shannon capacity is
\[
 \Theta(C_7)=\lim_{k\to\infty}
       \alpha(C_7^{\boxtimes k})^{1/k}
       =\sup_{k\ge1}\alpha(C_7^{\boxtimes k})^{1/k}.
\]
The limit exists because independent sets in two powers can be multiplied
to give an independent set in their combined power, making their
independence numbers supermultiplicative. All quantities are finite and
bounded after taking $k$th roots by seven.

Determine this real number exactly. In coding language, a set of tuples
is independent when any two distinct words have some coordinate whose
symbols are neither equal nor consecutive on the seven-cycle.
The capacity is the asymptotically achievable number of distinguishable
symbols per coordinate, not its logarithm.

The source discussion includes a factor $1/k$ inside its root; this
factor has limiting root one and does not alter the limit. The conventional
normalization above makes the supremum formula transparent. Computing one
finite strong power supplies a lower bound but does not by itself
determine the limit. The notation $\Theta$ distinguishes capacity from
other graph parameters sometimes denoted by a lowercase theta.
\subsection{Short English statement}
What is the exact exponential growth rate of the largest independent sets in repeated strong products of the seven-cycle?
\subsection{Research attempt}
\paragraph{Scope and process.}
GPT-6 Astra Ultra, 1 October 2026. This attempt keeps the catalogue statement above, checks primary literature, and develops a restricted argument with an adversarial gap audit. The proofs below are the mathematical output of this attempt, not a claim of novelty or external certification.

\paragraph{Literature and attack plan.}
The capacity remains a limit over every strong power; identifying one power does not identify it. Buys, Polak and Zuiddam's July 2026 primary preprint [R1] reports a substantially stronger lower bound, $3.258805369885\ldots$, than the elementary construction below. Its Lean verification is the authors' report; no Lean code was executed here. Our work independently determines the exact second-power independence number and audits the tensoring inference. It is a reproducible baseline, not a new best bound.

\paragraph{An explicit ten-word code.}
In $(\mathbb Z/7\mathbb Z)^2$, take
\[
\begin{split}
S=\{&(0,0),(1,2),(2,0),(2,4),(3,2),\\
     &(4,4),(4,6),(5,1),(6,3),(6,5)\}.
\end{split} \tag{1}
\]
Write its row sets as
\[
 A_0=\{0\},\ A_1=\{2\},\ A_2=\{0,4\},\ A_3=\{2\},\
 A_4=\{4,6\},\ A_5=\{1\},\ A_6=\{3,5\}.
\]
Every $A_i$ is independent in the seven-cycle. The union of any two consecutive rows, including rows $6,0$, is an independent set with no repeated column. Explicitly these unions are
\[
 \{0,2\},\ \{0,2,4\},\ \{0,2,4\},\ \{2,4,6\},\
 \{1,4,6\},\ \{1,3,5\},\ \{0,3,5\}.
\]
The cyclic gaps in every displayed set are at least two. Thus words in equal or consecutive rows are separated by a nonconfusable second coordinate. Words in other rows are already separated by their first coordinate. This proves that (1) is independent.

\paragraph{Matching upper bound for the second power.}
Let $T$ be any independent set in $C_7^{\boxtimes2}$ and let $B_i$ be its columns in row $i$. The sets $B_i,B_{i+1}$ are disjoint, and their union is independent in $C_7$: equality or adjacency of any two of those columns would make the corresponding words adjacent in the strong product. An independent set of $C_7$ has size at most three. Indeed, each selected vertex needs a distinct successor outside the set, so $2|B|\le7$.
Hence
\[
 |B_i|+|B_{i+1}|\le3\qquad(i\bmod7).
\]
Summing all seven inequalities gives $2|T|\le21$. Integrality yields $|T|\le10$, proving
\[
 \alpha(C_7^{\boxtimes2})=10. \tag{2}
\]

\paragraph{What tensoring certifies.}
Concatenating $t$ independently chosen words from $S$ gives an independent set $S^t$ in the $2t$th power: two different concatenations differ in some block, where (1) supplies a separating coordinate. Thus
\[
 \alpha(C_7^{\boxtimes2t})\ge10^t,
 \qquad \Theta(C_7)\ge\sqrt{10}.
\]
For comparison, a direct upper bound valid in every dimension is $\Theta(C_7)\le7/2$. To see this, let $H$ be an independent set in the $k$th power and let $Q=\{0,1\}^k$. Each translate $x+Q$ is a clique of the strong product, so it contains at most one element of $H$. There are $7^k$ translates and each element belongs to exactly $2^k$ of them. Double counting incidences gives $2^k|H|\le7^k$. The resulting interval is deliberately elementary and weaker than known literature bounds.

\paragraph{Verification and the first remaining gap.}
A Python check performed for this attempt compared all $\binom{10}{2}=45$ pairs in (1); for each pair at least one coordinate difference modulo seven was outside $\{0,1,6\}$. Its complete reproducible code is in \texttt{scripts/check\_attempt\_batch\_graphs\_2026\_10\_01.py}. The row argument separately proves optimality in dimension two, without trusting a search program. Cyclic adjacency between residues $0$ and $6$ was included in both checks.

Crucially, (2) does not imply $\alpha(C_7^{\boxtimes2t})=10^t$: a larger code need not factor into two-coordinate blocks. In fact [R1]'s lower bound exceeds $\sqrt{10}$, so promoting block optimality to capacity optimality would contradict checked research. The missing ingredient is a sharp bound for arbitrary dimension, matched by asymptotic constructions.

\paragraph{Final status.}
Exact second-power value proved, elementary general bounds proved, latest cited lower bound acknowledged. The exact Shannon capacity is \textbf{UNRESOLVED} by this attempt.

\subsection{Sources}
\begin{itemize}
\item[S1] ULAM.AI and the UnsolvedMath Contributors, supplied problems.json, record 3106 (OPG-36879), OpenGarden collection. Catalogue identity and source transcription; CC BY 4.0.
\url{https://huggingface.co/datasets/ulamai/UnsolvedMath}.
\item[S2] Open Problem Garden, Shannon capacity of the seven-cycle. Original problem and discussion; the mathematical formulation below is expanded from the supplied source record.
\url{http://www.openproblemgarden.org/op/shannon_capacity_of_the_seven_cycle}.
\item[R1] P. Buys, S. Polak and J. Zuiddam, \emph{Lean-verified lower bounds for the Shannon capacity of odd cycles}, arXiv:2607.29681 (31 July 2026), abstract checked 1 October 2026.
\url{https://arxiv.org/abs/2607.29681}.
\end{itemize}
\end{document}
